\scanpage{47}
Donc $\underline{\mathrm{Syst.lin.div}}_{X/S}$ est représentable ssi $\underline{\mathrm{Pic}}_{X/S}$ l'est.

\uline{Donnons nous} un syst. linéaire paramétré par $\check{P}/S$ de Br. Sév., et considérons la section correspondante $\xi$ de $\underline{\mathrm{Pic}}_{X/S}$.
\diagram{\begin{tikzcd}[column sep=large,row sep=large]
\underline L \arrow[d,phantom] & \underline L\otimes_S O_{\check P}(1)\otimes\underline N \arrow[d,phantom] \\
X \arrow[r,no head]\arrow[d,no head] & X\times_S\check P\arrow[d,no head] \\
S \arrow[r,no head] & \check P
\end{tikzcd}}

Je dis que $\xi$ \uline{est représentable par un $\underline L$ inversible sur $X$} \uline{ssi} \uline{$\check P$ est un fibré projectif}.
\margin{[$f:X\to S$ propre, $f_*(O_X)=O_S$ \unsure{univ\textsuperscript{t}}]}

Car la condition est suffisante \inserted{(\unsure{on sait par construction} $\unsure{C|\mathcal L}$ de $\xi$)}\figure{\begin{tikzpicture}\draw (0,0)..controls(-.6,-.1)and(-.6,-.5)..(-.6,-.8);\draw (3.5,-.1)..controls(4,-.6)and(3.8,-.9)..(3.4,-.9)--(1,-.9)..controls(.4,-.9)and(.1,-1.3)..(0,-1.7);\end{tikzpicture}}, prouvons qu'elle est nécessaire. \deleted{\unsure{Alors}} En effet, on a alors un Module de prés. finie $Q$ sur $S$, et un \inserted{mono}morphisme $\check P\longrightarrow P(Q)$, je dis que c'est un isom sur un $P(E)$, où $E$ est un quotient loc. libre de $Q$. Ce est par descente, \unsure{tjs|tjrs} avec le $=$ \unsure{principe} ---

\margin{Il suffit\par
$f_*(O_X)\simeq O_S$\par
et $f$ \deleted{\unsure{plat}} \inserted{de prés. finie}}

\margin{Si $X/S$ \unsure{adm} une section il est trivial que $\check P$ est ``trivial'', \unsure{car on sait qu'il|car alors qu'il} \unsure{devient trivial} sur $X$ .....}
\figure{\begin{tikzpicture}\draw (0,2.6) -- (0,0.3) .. controls (0,0.05) and (0.12,0) .. (0.5,0);\end{tikzpicture}}

\uline{Remarque} Partant d'un Br. Sév. \unsure{non trivial}, on construit un $\xi$ ne provenant pas d'un $\underline L$ \unsure{au $=$|avec} \unsure{la relation d'incidence can.|relation d'incidence can.}
\mspage{46}
